% =========================================================
% Chapter 5 -- ARMA models
% Author : Aymen Ben Brik (aymen.benbrik@esprit.tn)
% =========================================================

\chapter{ARMA models}
\label{ch:arma}

\section*{Introduction}
\addcontentsline{toc}{section}{Introduction}

Chapter~\ref{ch:stationarity} introduced stationarity and the
sample ACF/PACF; Chapter~\ref{ch:nonstationarity} showed how to
turn a non-stationary series into a stationary one by differencing.
This chapter assembles the central modelling family for stationary
time series: the \emph{autoregressive moving-average} processes,
ARMA$(p, q)$. We give their precise definitions, state the
stationarity and invertibility conditions in terms of polynomial
roots, derive their ACF/PACF \emph{signatures}, and detail the
practical workflow --- the Box--Jenkins method --- for identification,
estimation, diagnostics and forecasting.

% =========================================================
\section{Autoregressive process AR$(p)$}
\label{sec:ar}

\begin{defbox}[Autoregressive process AR$(p)$]
\label{def:ar}
A process $(X_t)_{t \in \Z}$ is autoregressive of order $p$ if
\[
  X_t = \phi_1 X_{t-1} + \phi_2 X_{t-2} + \dots + \phi_p X_{t-p}
        + \varepsilon_t,
  \qquad \varepsilon_t \sim \WN(0, \sigma^2),
\]
with $\phi_p \neq 0$. Equivalently, with $\Phi(z) = 1 - \phi_1 z -
\dots - \phi_p z^p$,
\[
  \Phi(B)\, X_t = \varepsilon_t.
\]
\end{defbox}

\begin{propbox}[Stationarity of AR$(p)$]
\label{prop:ar-stationary}
Equation $\Phi(B)\, X_t = \varepsilon_t$ admits a unique causal
stationary solution if and only if every root $z_i$ of
$\Phi(z) = 0$ satisfies $\abs{z_i} > 1$. In that case
\[
  X_t = \Phi^{-1}(B)\, \varepsilon_t
      = \sum_{j = 0}^{\infty} \psi_j\, \varepsilon_{t-j},
\]
with $\sum_{j} \abs{\psi_j} < \infty$.
\end{propbox}

\begin{exbox}[AR$(1)$ moments]
\label{ex:ar1}
For $X_t = \phi X_{t-1} + \varepsilon_t$ with $\abs\phi < 1$,
iterating gives
$X_t = \sum_{j \geq 0} \phi^j\, \varepsilon_{t - j}$, hence
\[
  \E[X_t] = 0,
  \qquad
  \Var(X_t) = \frac{\sigma^2}{1 - \phi^2},
  \qquad
  \rho(h) = \phi^{\abs h}.
\]
The ACF \textbf{decays geometrically} and the PACF \textbf{cuts off}
after lag~$1$ ($\alpha(1) = \phi$, $\alpha(h) = 0$ for $h \geq 2$).
\end{exbox}

\begin{propbox}[Yule--Walker equations]
\label{prop:yw}
Let $X_t \sim \AR{p}$ be stationary. Multiplying the equation by
$X_{t - h}$ and taking expectations gives, for $h = 1, \dots, p$,
\[
  \rho(h) = \phi_1\, \rho(h - 1) + \phi_2\, \rho(h - 2)
          + \dots + \phi_p\, \rho(h - p).
\]
In matrix form $\bm R \bm \phi = \bm r$, the system is invertible
under the stationarity condition; replacing $\rho$ by the sample
ACF $\widehat\rho$ gives the Yule--Walker estimator
$\widehat\phi_{\text{YW}}$.
\end{propbox}

\begin{exbox}[AR$(2)$ stationarity triangle]
For $X_t = \phi_1 X_{t-1} + \phi_2 X_{t-2} + \varepsilon_t$, the
roots of $1 - \phi_1 z - \phi_2 z^2 = 0$ lie outside the unit disk
iff
\[
  \phi_1 + \phi_2 < 1,
  \qquad
  \phi_2 - \phi_1 < 1,
  \qquad
  -1 < \phi_2 < 1.
\]
This is the Schur stability condition for the AR(2) polynomial.
\end{exbox}

% =========================================================
\section{Moving-average process MA$(q)$}
\label{sec:ma}

\begin{defbox}[Moving-average process MA$(q)$]
\label{def:ma}
A process $(X_t)_{t \in \Z}$ is moving-average of order $q$ if
\[
  X_t = \varepsilon_t + \theta_1 \varepsilon_{t-1}
        + \dots + \theta_q \varepsilon_{t-q},
  \qquad \varepsilon_t \sim \WN(0, \sigma^2),
\]
with $\theta_q \neq 0$. Equivalently, $X_t = \Theta(B)\,
\varepsilon_t$ with $\Theta(z) = 1 + \theta_1 z + \dots + \theta_q
z^q$.
\end{defbox}

\begin{propbox}[Moments and ACF of MA$(q)$]
\label{prop:ma-acf}
Every MA$(q)$ is (weakly) stationary, with
\[
  \E[X_t] = 0,
  \qquad
  \Var(X_t) = \sigma^2 (1 + \theta_1^2 + \dots + \theta_q^2),
\]
and autocovariance
\[
  \gamma(h) =
  \begin{cases}
    \sigma^2 \bigl(\theta_h + \theta_1 \theta_{h+1} + \dots
                 + \theta_{q-h} \theta_q\bigr)
                                & 1 \leq h \leq q, \\[4pt]
    0                           & h > q,
  \end{cases}
\]
with $\theta_0 = 1$. In particular $\rho(h) = 0$ for $h > q$:
the ACF \textbf{cuts off} at lag~$q$.
\end{propbox}

\begin{defbox}[Invertibility]
\label{def:invert}
An MA$(q)$ is \emph{invertible} if it admits an AR$(\infty)$
representation
\[
  \varepsilon_t = X_t - \pi_1 X_{t-1} - \pi_2 X_{t-2} - \dots,
  \qquad \sum_{j} \abs{\pi_j} < \infty.
\]
This holds if and only if every root of $\Theta(z) = 0$ lies
\emph{outside} the closed unit disk.
\end{defbox}

\begin{rembox}[Why invertibility matters]
Without invertibility, two parameter sets give the same ACF and so
are indistinguishable from the data. The MA$(1)$ models
$X_t = \varepsilon_t + \theta \varepsilon_{t-1}$ and
$X_t = \tilde\varepsilon_t + (1/\theta)\, \tilde\varepsilon_{t-1}$
have identical autocovariance functions; we conventionally choose
the invertible one, $\abs\theta < 1$.
\end{rembox}

% =========================================================
\section{ARMA$(p, q)$}
\label{sec:arma}

\begin{defbox}[ARMA$(p, q)$]
\label{def:arma}
A process $(X_t)$ is ARMA$(p, q)$ if it satisfies
\[
  \Phi(B)\, X_t = \Theta(B)\, \varepsilon_t,
  \qquad \varepsilon_t \sim \WN(0, \sigma^2),
\]
with $\Phi(z) = 1 - \phi_1 z - \dots - \phi_p z^p$,
$\Theta(z) = 1 + \theta_1 z + \dots + \theta_q z^q$, and $\Phi$,
$\Theta$ sharing no common root.
\end{defbox}

\begin{propbox}[Stationarity, invertibility and existence]
\label{prop:arma}
The equation $\Phi(B) X_t = \Theta(B) \varepsilon_t$ admits a
unique causal stationary invertible solution if and only if
\[
  \abs{z_i} > 1 \text{ for every root of } \Phi,
  \qquad
  \abs{z_j} > 1 \text{ for every root of } \Theta.
\]
The solution is then
\[
  X_t = \Phi^{-1}(B)\, \Theta(B)\, \varepsilon_t
      = \sum_{j = 0}^{\infty} \psi_j\, \varepsilon_{t-j}.
\]
\end{propbox}

\begin{rembox}
The MA$(\infty)$ coefficients $(\psi_j)$ are obtained by polynomial
division of $\Theta(z)$ by $\Phi(z)$ around $z = 0$. They
control both the autocovariance and the forecast variance
(Section~\ref{sec:forecasting}).
\end{rembox}

% =========================================================
\section{Box--Jenkins identification}
\label{sec:identification}

The choice of $(p, q)$ is guided by the sample ACF and PACF, whose
\emph{signatures} differ across model classes:

\begin{center}
\renewcommand{\arraystretch}{1.4}
\begin{tabular}{lll}
\toprule
Model        & ACF $\rho(h)$              & PACF $\alpha(h)$ \\
\midrule
AR$(p)$      & exponential decay          & cuts off after lag $p$ \\
MA$(q)$      & cuts off after lag $q$     & exponential decay \\
ARMA$(p, q)$ & exponential decay (after $q$) & exponential decay (after $p$) \\
\bottomrule
\end{tabular}
\end{center}

\begin{rembox}[Bartlett bands]
``Cuts off'' is read as $\widehat\rho(h)$ falling \emph{inside} the
$\pm 1.96 / \sqrt T$ band for $h$ larger than the cutoff lag.
``Decays'' means the values stay outside the band but shrink in
absolute value.
\end{rembox}

\begin{exbox}[Reading a sample ACF/PACF]
On a series of length $T = 200$, the sample PACF shows
$\widehat\alpha(1) = 0.62$, $\widehat\alpha(2) = -0.04$,
$\widehat\alpha(3) = 0.02$ and the sample ACF decays geometrically.
The Bartlett band is $\pm 1.96 / \sqrt{200} \approx \pm 0.139$, so
$\widehat\alpha(h)$ is inside the band for $h \geq 2$:
the PACF cuts off at lag~$1$, suggesting AR$(1)$.
\end{exbox}

% =========================================================
\section{Estimation}
\label{sec:arma-estimation}

\subsection{Yule--Walker (AR only)}

For pure AR$(p)$, plug $\widehat\rho$ into
Proposition~\ref{prop:yw} to obtain the Yule--Walker estimate.
This is method-of-moments: \emph{consistent} but not asymptotically
efficient.

\subsection{Conditional sum of squares (CSS)}

Conditioning on the first $p$ observations and assuming
$\varepsilon_t = 0$ for $t \leq 0$, the CSS estimator is
\[
  \widehat{\bm\theta}_{\text{CSS}}
  = \arg\min_{\bm\theta} \sum_{t = p + 1}^{T}
        \widehat\varepsilon_t(\bm\theta)^2,
\]
where $\widehat\varepsilon_t$ is computed recursively from the
ARMA equation. CSS is fast and equivalent to MLE in large samples
under Gaussian innovations.

\subsection{Maximum likelihood (MLE)}

Under $\varepsilon_t \iid \mathcal N(0, \sigma^2)$, the joint
distribution of $(X_1, \dots, X_T)$ is multivariate normal with
covariance matrix $\Gamma_T(\bm\theta)$. The log-likelihood
\[
  \log \mathcal L(\bm\theta, \sigma^2)
  = -\tfrac{T}{2}\log(2\pi)
    -\tfrac{1}{2}\log \det \Gamma_T(\bm\theta)
    -\tfrac{1}{2} \bm X^\top \Gamma_T(\bm\theta)^{-1} \bm X
\]
is maximised numerically. The standard implementation in
\texttt{statsmodels} uses a state-space representation and the
Kalman filter, which scale linearly in $T$ and provide asymptotic
standard errors.

\begin{thbox}[Asymptotic distribution of the MLE]
Under regularity assumptions (stationarity, invertibility, no
common roots), as $T \to \infty$,
\[
  \sqrt{T}\, (\widehat{\bm\theta}_{\text{MLE}} - \bm\theta_0)
  \;\xrightarrow{d}\; \mathcal N(0, I(\bm\theta_0)^{-1}),
\]
where $I(\bm\theta_0)$ is the Fisher information matrix evaluated
at the true parameter $\bm\theta_0$.
\end{thbox}

% =========================================================
\section{Model selection}
\label{sec:selection}

\begin{defbox}[Information criteria]
\label{def:aic-bic}
For a fitted ARMA$(p, q)$ with $k = p + q + 1$ free parameters
(including $\sigma^2$),
\[
  \mathrm{AIC} = -2 \log \widehat{\mathcal L} + 2 k,
  \qquad
  \mathrm{BIC} = -2 \log \widehat{\mathcal L} + k \log T.
\]
The chosen order minimises the criterion.
\end{defbox}

\begin{rembox}[AIC vs BIC]
AIC favours predictive accuracy and tends to over-fit on small
samples; BIC penalises complexity more and is consistent for
the true order under standard assumptions. When the two disagree,
prefer the BIC choice for parsimony and confirm with the residual
diagnostics.
\end{rembox}

% =========================================================
\section{Residual diagnostics}
\label{sec:diagnostics}

\begin{rembox}[Adequacy checks]
A well-specified ARMA model produces residuals
$\widehat\varepsilon_t = X_t - \widehat X_{t \mid t-1}$ that are
\emph{indistinguishable from white noise}. Standard checks:
\begin{itemize}
  \item time plot of $\widehat\varepsilon_t$ -- mean stable around
        zero, variance constant;
  \item sample ACF of $\widehat\varepsilon_t$ -- spikes inside
        Bartlett bands;
  \item Ljung--Box test;
  \item QQ-plot vs $\mathcal N(0, 1)$ for normality (Jarque--Bera
        test).
\end{itemize}
\end{rembox}

\begin{thbox}[Ljung--Box test on residuals]
\label{th:lb-residuals}
\[
  Q_{\text{LB}}(m)
  = T(T + 2) \sum_{h = 1}^{m}
        \frac{\widehat\rho_{\widehat\varepsilon}(h)^2}{T - h}.
\]
Under the null that $\widehat\varepsilon_t$ is white noise, after
fitting an ARMA$(p, q)$,
\[
  Q_{\text{LB}}(m) \;\xrightarrow{d}\; \chi^2_{m - p - q}.
\]
Reject the null if $Q_{\text{LB}}(m)$ exceeds the
$\chi^2_{m - p - q, 1 - \alpha}$ quantile.
\end{thbox}

% =========================================================
\section{Forecasting}
\label{sec:forecasting}

\begin{defbox}[Best linear forecast]
\label{def:forecast}
The minimum mean-squared-error linear forecast at horizon $h$,
given the past $\mathcal F_T = \sigma(X_T, X_{T-1}, \dots)$, is
\[
  \widehat X_{T + h \mid T}
  = \E\bigl[X_{T + h} \mid \mathcal F_T \bigr].
\]
The forecast error is
$e_{T + h} = X_{T + h} - \widehat X_{T + h \mid T}$.
\end{defbox}

\begin{propbox}[Forecast error variance]
\label{prop:forecast-var}
With the MA$(\infty)$ representation
$X_t = \sum_{j \geq 0} \psi_j\, \varepsilon_{t - j}$,
\[
  e_{T + h} = \sum_{j = 0}^{h - 1} \psi_j\, \varepsilon_{T + h - j},
\]
hence
\[
  \E\bigl[e_{T + h}\bigr] = 0,
  \qquad
  \Var\bigl(e_{T + h}\bigr)
  = \sigma^2 \sum_{j = 0}^{h - 1} \psi_j^2.
\]
\end{propbox}

\begin{exbox}[Forecast for AR$(1)$]
With $X_t = \phi X_{t-1} + \varepsilon_t$, $\abs\phi < 1$:
\[
  \widehat X_{T + h \mid T} = \phi^h X_T,
  \qquad
  \Var(e_{T + h}) = \sigma^2 \frac{1 - \phi^{2h}}{1 - \phi^2}.
\]
As $h \to \infty$, $\widehat X_{T + h \mid T} \to 0$ (the
unconditional mean) and the variance tends to
$\sigma^2 / (1 - \phi^2)$ (the unconditional variance).
\end{exbox}

\begin{rembox}[Confidence intervals]
Under Gaussian innovations, the $1 - \alpha$ confidence interval is
\[
  \widehat X_{T + h \mid T}
  \;\pm\; z_{1 - \alpha / 2}\, \sigma\,
        \sqrt{\sum_{j = 0}^{h - 1} \psi_j^2}.
\]
For stationary ARMA, the width converges to a finite limit; for
ARIMA (Chapter~\ref{ch:nonstationarity}), it grows without bound.
\end{rembox}

% =========================================================
\section*{Summary}
\addcontentsline{toc}{section}{Summary}

\noindent\textbf{Key takeaways.}
\begin{itemize}
  \item AR$(p)$ depends on past values; MA$(q)$ depends on past
        shocks; ARMA$(p, q)$ combines both.
  \item Stationarity = roots of $\Phi$ outside the unit disk;
        invertibility = roots of $\Theta$ outside the unit disk.
  \item ACF/PACF signatures: AR cuts the PACF, MA cuts the ACF,
        ARMA shows double exponential decay.
  \item Estimation: Yule--Walker for pure AR; CSS or MLE for
        general ARMA. \texttt{statsmodels.tsa.arima.ARIMA} provides
        both.
  \item Model selection by AIC / BIC; validate with Ljung--Box on
        residuals.
  \item Forecasts widen with horizon according to
        $\Var(e_{T + h}) = \sigma^2 \sum_{j = 0}^{h - 1} \psi_j^2$.
\end{itemize}
